\section{Useful Memory Geometry}\label{sec:memory_geometry}

Under continuous drift, memory is not only a statistical resource but also a
temporal liability. The core question is therefore not whether one can react to
change, but how long a sample should remain operationally useful before it
becomes stale.

\begin{theorem}[Lag-Inclusive Tracking Error Bound]\label{thm:lag_error}
Let $X_{t-j} \sim P^*_{t-j}$ and let
\[
  \widehat P_t^{(n)} \,=\, \frac{1}{n}\sum_{j=0}^{n-1} \delta_{X_{t-j}}
\]
be the uniform-window empirical estimator. Assume the target path is locally
$W_2$-Lipschitz with rate $\zeta$, so that
\[
  W_2\!\left(P^*_{t-j}, P^*_t\right) \leq \zeta j
  \qquad \text{for all } 0 \leq j < n,
\]
and assume the finite-sample term obeys
\[
  \mathbb{E}\!\left[W_2\!\left(\widehat P_t^{(n)}, \bar P_t^{(n)}\right)\right]
  \leq C_K n^{-1/2},
\]
where $\bar P_t^{(n)} := n^{-1}\sum_{j=0}^{n-1} P^*_{t-j}$. Then the expected
tracking error satisfies
\[
  \mathcal{E}(n) \;\leq\; C_K\,n^{-1/2} \;+
  \tfrac{1}{2}\,\zeta\,n,
\]
where the first term is the finite-sample statistical cost and the second term
is the staleness cost induced by retaining old information.
\end{theorem}

\begin{lemma}[Finite-sample Wasserstein term]\label{lem:finite_sample_term}
Let $Q$ be a probability measure on $\mathbb R^{d_\mathrm{eff}}$ with support
contained in a ball of radius $R$ and finite second moment. Let
$X_1,\dots,X_n$ be i.i.d. samples from $Q$, and let
$\hat Q_n := n^{-1}\sum_{i=1}^n \delta_{X_i}$. Then:
\begin{itemize}
  \item If $d_\mathrm{eff}<4$, there exists a constant $C_{d_\mathrm{eff},R}$
  such that
  \[
    \mathbb E\,W_2(\hat Q_n,Q) \le C_{d_\mathrm{eff},R}\,n^{-1/2}.
  \]
  \item If $d_\mathrm{eff}=4$, the same bound holds with an additional
  $(\log n)^{1/2}$ factor.
  \item If $d_\mathrm{eff}>4$, the raw Wasserstein rate is slower than
  $n^{-1/2}$, so the constant-only form above should not be used.
\end{itemize}
If the geometry is computed with debiased Sinkhorn divergence $S_\varepsilon$
on a bounded domain of diameter $D$, then for fixed $\varepsilon>0$ the
empirical fluctuations remain $O(n^{-1/2})$ with a constant
$L(D,\varepsilon)$ that is dimension-free at fixed $\varepsilon$ and increases
as $\varepsilon\downarrow 0$.

A short numerical sanity check for this regime is reported in
Appendix~\ref{app:finite_sample_geometry}; it is not used as a theorem, but it
matches the intended low-dimensional interpretation of the finite-sample term.
\end{lemma}

The two terms in \cref{thm:lag_error} pull in opposite directions. Increasing
$n$ reduces variance but raises the age of the retained evidence. That tension is
the basic coherence-lag trade-off that governs useful memory under drift.

\begin{proposition}[Optimal Window Under Drift]\label{prop:coherence_delay}
Under \cref{thm:lag_error}, the error $\mathcal{E}(n)$ is minimised at
\[
  n^* \;=\; \left(\frac{C_K}{\zeta}\right)^{2/3},
\]
with minimum value
\[
  \mathcal{E}_{\min} \;=\; \tfrac{3}{2}\,C_K^{2/3}\,\zeta^{1/3}.
\]
\end{proposition}

\begin{proof}
Apply the triangle inequality for $W_2$:
\[
  W_2\!\left(\widehat P_t^{(n)}, P_t^*\right)
  \leq
  W_2\!\left(\widehat P_t^{(n)}, \bar P_t^{(n)}\right)
  + W_2\!\left(\bar P_t^{(n)}, P_t^*\right).
\]
The first term is controlled by the assumed finite-sample bound.
For the second term, let $\pi_j$ be an optimal coupling between
$P^*_{t-j}$ and $P_t^*$. Then the averaged coupling
$\bar\pi := n^{-1}\sum_{j=0}^{n-1} \pi_j$ couples $\bar P_t^{(n)}$ with
$P_t^*$, so the transport cost is at most the average of the individual costs.
Combined with the $W_2$-Lipschitz drift bound, this gives
\[
  W_2\!\left(\bar P_t^{(n)}, P_t^*\right)
  \leq \frac{1}{n}\sum_{j=0}^{n-1} W_2\!\left(P^*_{t-j}, P_t^*\right)
  \leq \frac{1}{n}\sum_{j=0}^{n-1} \zeta j
  \leq \tfrac{1}{2}\,\zeta n.
\]
Combining the two bounds yields the stated inequality.
\end{proof}

\begin{remark}[Measurement layer]
The paper's theorem is stated in $W_2$, so $\varepsilon$ does not enter the
mathematical drift law itself. In the synthetic measurement layer, however, the
observable geometry is often a debiased Sinkhorn divergence with fixed
$\varepsilon$, and then the effective finite-sample constant should be read as
$C_K^{\mathrm{eff}} = C_{d_\mathrm{eff},R} + L(D,\varepsilon)$: the raw
Wasserstein part depends on intrinsic dimension, while the Sinkhorn part is
dimension-free at fixed $\varepsilon$ but worsens as $\varepsilon\downarrow 0$.

In the synthetic benchmarks, $\hat\zeta_t$ is estimated from short-window
block differences using an EMA update, $C_K$ is calibrated from a stationary
prefix, and $\varepsilon$ is held fixed within each experiment family, with
runtime sensitivity to $\varepsilon$ reported separately. The benchmark-specific
drift proxy, robustness check, and sensitivity sweep are detailed in
Appendix~\ref{app:zeta}.
The same appendix also validates an observable prequential calibration rule for
choosing $\alpha$ and the block scale, so the online regulator is not tied to a
fixed manual setting.
For the horizon constant itself, the prefix calibration is the simple stationary
baseline used to estimate the scale in the synthetic stream; it is not a tuned
oracle, but it is also not arbitrary.

The Sinkhorn-to-$W_2$ gap is therefore treated as an empirical calibration
constant. In the tested Gaussian shift sweep, the high-dimensional slice shows
that this gap is nearly flat in sample size for small $\varepsilon$ and decays
with $n$ only when $\varepsilon$ is large; the resulting horizon inflation in the
sampled grid ranges from about $1.36\times$ to $4.24\times$. This supports the
paper's use of fixed-$\varepsilon$ Sinkhorn as a measurement layer while keeping
the theoretical law in $W_2$. For fixed $\varepsilon$ on bounded domains, the
debiased Sinkhorn divergence converges at rate $n^{-1/2}$ with
$\varepsilon$-dependent constants~\citep{GenevayEtAl2019}; the approximation gap
to $W_2^2$ is $O(\varepsilon)$ and vanishes as $\varepsilon\to 0$. We therefore
treat the Sinkhorn layer as a calibration proxy rather than part of the law
itself.
\end{remark}

\begin{corollary}[EWMA Analogue]\label{cor:ewma_cube_root}
For exponential weights $w_j = \alpha(1-\alpha)^j$, the effective lag is
$\tau_\alpha = (1-\alpha)/\alpha$ and the effective sample size scales as
$n_{\mathrm{eff}} \asymp 1/\alpha$. Balancing the same statistical and
staleness terms therefore yields the same cube-root exponent in the effective
memory scale.
\end{corollary}

\begin{definition}[Effective Memory Age]\label{def:memory_age}
Let an estimator place weights $w_0,\dots,w_{n-1}$ on the most recent
observations, with $\sum_j w_j = 1$. Its \emph{effective memory age} is
\[
  A(w) \;=\; \sum_{j=0}^{n-1} j\,w_j,
\]
and its corresponding freshness proxy is the monotone transform
$F(w) = 1/(1+A(w))$. The role of $F(w)$ is purely geometric: it records how old
the retained evidence is before any local drift scale is specified. For a
uniform window, $A(w)=(n-1)/2$.
\end{definition}

\begin{definition}[Temporal Coherence]\label{def:temporal_coherence}
For an estimator with weights $w$ at time $t$, define its \emph{operational
staleness}
\[
  S_t(w) \;=\; \sum_{j=0}^{n-1} w_j\,
  W_2\!\left(\Pot^*(S,t-j),\Pot^*(S,t)\right).
\]
Its \emph{temporal coherence} is the inverse alignment score
\[
  C_t(w) \;=\; \frac{1}{1 + S_t(w)}.
\]
If the target path is locally $W_2$-Lipschitz at rate $\zeta_t$, in the sense
that $W_2\!\left(\Pot^*(S,t-j),\Pot^*(S,t)\right) \leq \zeta_t j$ for the
relevant lags, then
\[
  S_t(w) \;\leq\; \zeta_t A(w),
  \qquad
  C_t(w) \;\geq\; \frac{1}{1+\zeta_t A(w)}.
\]
In that sense, $F(w)$ is the rate-free freshness proxy, while $C_t(w)$ measures
operational alignment with the current distribution.
\end{definition}

The quantity $A(w)$ gives a compact description of the temporal geometry of a
finite-memory estimator. Short windows preserve temporal coherence at the cost
of variance; long windows reduce variance at the cost of staleness. The
cube-root horizon marks the point where that exchange stops improving. For a
uniform window under
approximately constant local drift, the induced staleness scales as
$S_t(w)\approx \zeta(n-1)/2$, which is the same age term that appears in
\cref{thm:lag_error} up to constants.

\begin{proposition}[Finite-Memory Floor via Gaussian Location Subclass]\label{prop:window_minimax}
Let $\mathcal P_{\zeta,m}$ be the class of length-$m$ distribution paths for
which the following hold:
\begin{enumerate}
  \item the path is $W_2$-Lipschitz with rate at most $\zeta$;
  \item the lower-bound construction is carried out on the Gaussian location
  subclass $X_t \sim \mathcal N(\mu_t,\sigma^2 I_d)$ with
  $\|\mu_t-\mu_s\|_2 \le \zeta |t-s|$;
  \item the estimator at time $m$ only uses the last $m$ samples.
\end{enumerate}
Measure risk by the endpoint $W_2$ error, which coincides with Euclidean
endpoint error in the Gaussian location subclass. Then, in the critical-window
regime, the minimax tracking error over that class cannot beat the same
variance-plus-staleness scaling:
\[
  \inf_{\hat\Pot_m}\sup_{\Pot^*}\,\mathbb{E}\,d\bigl(\hat\Pot_m,\Pot^*\bigr)
  \;\gtrsim\; \max\!\left\{C_K\,m^{-1/2},\; c\,\sigma^{2/3}\zeta^{1/3}\right\}.
\]
Consequently, the cube-root horizon is not merely the optimum of one estimator;
it is the structural floor of finite memory on this drift class.
\end{proposition}

\paragraph{Proof.}
Write $R_m^* \coloneqq \inf_{\delta_m} R_m(\delta_m)$. We exhibit two
independent obstructions. The first is the ordinary finite-sample estimation
floor, obtained by restricting to the constant-path subclass $\mu_s \equiv
\theta$, which is contained in the admissible drift class. For the two
hypotheses $\theta_\pm = \pm \sigma /(2\sqrt{m})$, the induced product-Gaussian
laws have Kullback--Leibler divergence
\[
  \mathrm{KL}(P_+,P_-)
  = \frac{1}{2\sigma^2} \sum_{j=1}^{m} (\theta_+ - \theta_-)^2
  = \frac{1}{2\sigma^2} m\left(\frac{\sigma}{\sqrt{m}}\right)^2
  = \frac{1}{2}.
\]
For any estimator $\delta_m$, the plug-in test $\psi = \mathbf{1}\{\delta_m \geq 0\}$
must err whenever the estimate falls on the wrong side of the origin; on that
event the endpoint error is at least $\sigma /(2\sqrt{m})$. By Le Cam's
two-point lemma and Pinsker's inequality, the average testing error is at least
\[
  \frac{1-\mathrm{TV}(P_+,P_-)}{2}
  \geq
  \frac{1-\sqrt{\mathrm{KL}(P_+,P_-)/2}}{2}
  \geq \frac{1}{4}.
\]
Hence
\[
  R_m^* \;\geq\; \frac{\sigma}{8\sqrt{m}},
\]
so one valid explicit choice is $c_0 = 1/8$.

The second obstruction is local drift. Fix any $1\leq h\leq m$ and define
\[
  \beta_h \coloneqq \min\!\left\{\zeta,\; \frac{\sigma}{4h^{3/2}}\right\},
  \qquad
  \mu^{\pm,h}_{t-j} \coloneqq \pm \,\beta_h\,(h-j)_+, 
  \qquad 0\leq j\leq m-1,
\]
where $(u)_+ \coloneqq \max\{u,0\}$. Each path belongs to the admissible
class because its slope magnitude is at most $\beta_h \leq \zeta$. The
endpoint gap is
\[
  |\mu_t^{+,h} - \mu_t^{-,h}| = 2\beta_h h.
\]
The corresponding Gaussian product laws satisfy
\[
  \mathrm{KL}(P_{+,h},P_{-,h})
  = \frac{1}{2\sigma^2} \sum_{j=0}^{m-1}
    \bigl(2\beta_h (h-j)_+\bigr)^2
  \leq \frac{2\beta_h^2}{\sigma^2} \sum_{r=1}^{h} r^2
  \leq \frac{2\beta_h^2}{\sigma^2} h^3
  \leq \frac{1}{8}.
\]
Applying the same reduction from estimation to testing, now with endpoint error
threshold $\beta_h h$, and using Le Cam together with Pinsker gives average
testing error at least
\[
  \frac{1-\mathrm{TV}(P_{+,h},P_{-,h})}{2}
  \geq
  \frac{1-\sqrt{\mathrm{KL}(P_{+,h},P_{-,h})/2}}{2}
  \geq \frac{3}{8}.
\]
Therefore
\[
  R_m^* \;\geq\; \frac{3}{8}\,\beta_h h
  \;\geq\;
  \frac{3}{32}\,\min\!\left\{\zeta h,\; \sigma h^{-1/2}\right\}
\]
which is again a fully explicit numerical lower bound. Since $h$ was arbitrary,
\[
  R_m^* \;\geq\;
  \frac{3}{32}\,\sup_{1\leq h\leq m}
  \min\!\left\{\zeta h,\; \sigma h^{-1/2}\right\}.
\]

Taking the maximum of the constant-path and ramp bounds proves the displayed
lower bound. For the cube-root regime, assume
$1 \leq h_\star \coloneqq (\sigma/\zeta)^{2/3} \leq m$. Choose an integer
$\bar h \in [h_\star/2, h_\star]$. Then
\[
  \zeta \bar h \geq \frac{1}{2}\,\zeta h_\star
  = \frac{1}{2}\,\sigma^{2/3}\zeta^{1/3},
\]
and
\[
  \sigma \bar h^{-1/2}
  \geq \sigma h_\star^{-1/2}
  = \sigma^{2/3}\zeta^{1/3}.
\]
Therefore
\[
  \min\!\left\{\sigma \bar h^{-1/2},\; \zeta \bar h\right\}
  \geq \frac{1}{2}\,\sigma^{2/3}\zeta^{1/3},
\]
and the ramp bound gives
\[
  R_m^* \geq \frac{3}{64}\,\sigma^{2/3}\zeta^{1/3}.
\]
So one admissible explicit constant in the critical-window regime is
$c = 3/64$.

\begin{remark}[Scope of the lower bound]\label{rem:lower_bound_scope}
The lower bound is parametric only in the testing construction: it uses a
Gaussian location subclass whose mean path is itself $W_2$-Lipschitz, and the
distance in the lower bound is the endpoint $W_2$ error (Euclidean on that
subclass). That is enough to make the result minimax for any larger class
containing that subclass, including the paper's general $W_2$-Lipschitz path
class. So the proposition is not claiming that every path is Gaussian; it is
claiming that the worst case already appears inside the Gaussian mean-shift
subproblem.
Equivalently, since the Gaussian location family is a subfamily of the ambient
$W_2$-Lipschitz class at the same drift rate $\zeta$, any lower bound proven on
that subclass is automatically a lower bound for the larger class.
\end{remark}

\begin{remark}[Measurement layer]
Debiased Sinkhorn divergence serves as a practical measurement layer when the
underlying geometry must be estimated from samples. That computational choice
supports the tracking law, but the law itself is about the horizon of useful
memory, not about any particular OT solver.
\end{remark}
